% s-131-p131-m002; original PDF 131
\Aphorism{s-131-p131-m002}{181}
Die Spielereien und Äusserlichkeiten der Liebe, Begegnungen, Briefe, Bilder, haben ihren guten Sinn: sie halten das geliebte Object über der Bewusstseinsschwelle des Liebenden und verhindern, dass es ins Unterbewusste sinkend eine tiefe Störung aller Functionen anrichte. Diese Hygiene darf man nicht versäumen; lässt man die Liebe erst ins Blut gehen, so ist eine fundamentale Verstimmung da, der „von oben“ nicht mehr beizukommen ist. Man soll das Animalische nicht ausschalten, weil sonst die organische Reaction ungemildert eintritt; man muss die wirkliche Geliebte als rettende Kühlung auflegen, um nicht in den Ekstasen der Phantasieliebe zu verbrennen. Das Bewusstsein, das Bilderbuch und Spielzeug des Kindes Mensch, wirkt hier als Tröstung, als Chinin gegen das erotische Fieber, als Ernüchterung von einem gefährlichen und abspannenden Rausch.
\par

